\documentclass[12pt,a4paper]{article}
\usepackage[utf8]{inputenc}
\usepackage[T1]{fontenc}
\usepackage{mathptmx}
\usepackage{amsmath,amssymb}
\usepackage{graphicx}
\usepackage[margin=1in]{geometry}
\usepackage{booktabs}
\usepackage{multirow}
\usepackage{placeins}
\usepackage{hyperref}
\usepackage{natbib}
\usepackage{caption}
\usepackage{float}
\usepackage{setspace}
\usepackage{xcolor}
\usepackage{xurl}
\onehalfspacing
\setlength{\emergencystretch}{3em}
\hypersetup{colorlinks=true,linkcolor=blue!60!black,citecolor=blue!60!black,urlcolor=blue!60!black,breaklinks=true}

\newcommand{\rbar}{\langle r\rangle}

\title{\textbf{Charge-and-Release Rhythms in Ore-Forming Systems Are\\
Scalar Clocks, Not Level Repulsion:\\
A Cautionary Random-Matrix Analysis of Metallogenic Timing}\\[4pt]
\large Version 2.0 (erratum). A single ore system is a relaxation\\
oscillator---a scalar renewal clock---structurally not a GOE spectrum}

\author{Ruqing Chen\\
\textit{GUT Geoservice Inc., Montreal, Canada}\\
\texttt{ruqing@hotmail.com}}
\date{June 2026}

\begin{document}
\maketitle

\begin{abstract}
\noindent\textbf{This is version~2.0, a self-correcting erratum.} Version~1 of
this work claimed that the multi-stage timing of single ore-forming systems
carries the Gaussian-orthogonal/unitary-ensemble (GOE/GUE) \emph{level-repulsion}
signature of random matrix theory (RMT), interpreted as a charge-and-release
``memory.'' That central claim was wrong, and we retract it here. Re-auditing the
identical data with the diagnostics the original analysis omitted---the shuffle
(permutation) test, the coefficient of variation (CV), an i.i.d.\ renewal-clock
comparison, and an age-error Monte Carlo---shows that the high spacing ratios are
a \emph{scalar-clock} (narrow-marginal renewal) signature, not sequential level
repulsion. The pooled spacing ratios are reproduced \emph{exactly}
($\rbar=0.60$ Andean, 0.57 Nanling, 0.71 Tethyan, 0.68 orogenic gold), but (i)~they
are invariant under shuffling, so they reflect the marginal spacing distribution,
not the sequence; (ii)~a zero-repulsion gamma renewal clock with the same CV
reproduces every value; (iii)~the small negative ``rigidity'' (obs$-$shuf
$=-0.031$ for gold) is fully reproduced by pushing a \emph{zero-repulsion} null
through the same within-deposit mean-normalization---it is a closed-data
(compositional) artifact, not spectral rigidity; and (iv)~within the quoted
geochronological errors the pooled $\rbar$ ranges across the entire
Poisson-to-GUE axis, so the point estimate is unresolvable at these sample sizes
($n=21$--$40$, far below the $n\gtrsim300$ needed to distinguish a scalar clock
from GOE)---and the $\pm1$--$5$~Ma analytical errors of the underlying chronometers
are themselves comparable to the intervals, so isotopic dating cannot in principle
resolve the sequential rigidity a Wigner--Dyson spectrum would require. A charge-and-release reservoir is a classical relaxation oscillator,
i.e.\ a renewal process; by the Bohigas--Giannoni--Schmit criterion it is
structurally \emph{not} a chaotic GOE spectrum. \textbf{What survives is the
superposition result:} pooling many independent sources drives the local
statistic toward Poisson ($\rbar\to0.39$), a genuine and retained finding. We
keep the four-province dataset and the superposition analysis, retract the
level-repulsion interpretation, and reframe single ore systems as scalar
charge-release clocks. Basis of the correction: the scalar-clock audit note,
Zenodo \url{https://doi.org/10.5281/zenodo.20980670}.

\medskip
\noindent\textbf{Keywords:} random matrix theory; scalar clock; renewal process;
relaxation oscillator; metallogeny; ore geochronology; spectral superposition;
self-correction

\medskip
\noindent\textbf{Code/data:} \url{https://github.com/Ruqing1963/metallogeny-rmt}
\end{abstract}

\section{Changes from version~1 (erratum)}
\label{sec:changes}

This section states plainly what changed and why, so that readers of version~1 can
see exactly which claims are withdrawn.

\begin{itemize}
\item \textbf{Retracted.} The claim that single ore-forming systems exhibit
GOE/GUE \emph{level repulsion}, and the associated ``intrinsic charge-release
memory'' and ``grand unification of five deep-Earth systems'' framing. The
non-magmatic orogenic-gold ``control'' does \emph{not} establish intrinsic
repulsion; like the magmatic provinces it is a scalar clock.
\item \textbf{Why.} The spacing ratio $\rbar$ is two-ended: a high value can come
either from genuine sequential level repulsion \emph{or} from a narrow marginal
(scalar) spacing distribution. Version~1 read a high $\rbar$ as repulsion without
running the shuffle test that separates the two. A charge-and-release reservoir
is a relaxation oscillator---a renewal process with a characteristic recurrence
time---whose narrow interval distribution produces a high $\rbar$ with \emph{no}
sequential correlation. By the Bohigas--Giannoni--Schmit criterion
\citep{bgs1984}, genuine GOE requires a chaotic, strongly-mixed spectrum, which an
integrable relaxation oscillator is not.
\item \textbf{Retained (real).} The four-province compiled dataset, and the
\emph{superposition} result: pooling intervals across many independent sources
drives the spacing ratio toward Poisson ($\rbar\to0.386$). This is the standard
superposition theorem for point processes \citep{daley2003} and is unaffected by
the correction.
\item \textbf{Added.} The shuffle (permutation) test on every series; the
CV of the pooled spacings and an i.i.d.\ gamma renewal-clock comparison; a
closed-data control showing the apparent negative ``rigidity'' is a normalization
artifact; and an age-error Monte Carlo showing the point estimate is
unresolvable. New code: \texttt{code/audit\_scalar\_clock.py}.
\item \textbf{Basis.} The methodological diagnoses are established in the
scalar-clock audit note, Zenodo \url{https://doi.org/10.5281/zenodo.20980670}.
\item \textbf{Preserved.} Version~1 (text, figures, results) is preserved in the
git history (tag \texttt{v1.0}), in \texttt{paper/archive/v1/}, and as the prior
Zenodo version. License unchanged (MIT).
\end{itemize}

\section{Introduction}

Ore deposits form in pulses. A single porphyry copper system, granite-related
W--Sn deposit, or orogenic gold camp typically records several discrete
mineralization stages \citep{sillitoe2010,chelle2017}. Version~1 of this paper
asked whether the timing of those pulses carries the level-repulsion signature of
random matrix theory \citep{wigner1957,mehta2004}, motivated by a
\emph{charge-and-release hypothesis}: each pulse depletes the local metal, fluid,
and heat budget, which must re-accumulate before the next, imposing a minimum
recurrence interval. Version~1 reported high spacing ratios in four provinces and
read them as GOE/GUE repulsion.

That reading conflated two physically distinct situations. The spacing ratio
$\rbar$ \citep{oganesyan2007,atas2013} is \emph{two-ended}. A value well above the
Poisson expectation ($0.386$) can arise from (i)~genuine sequential level
repulsion, in which the position of each level depends on its neighbours
(spectral rigidity), or (ii)~a narrow \emph{marginal} distribution of spacings
with \emph{no} sequential structure---a ``scalar clock.'' A relaxation oscillator
that fires at a roughly fixed recurrence time (a charge-and-release reservoir, a
dripping tap, a geyser) is exactly case~(ii): its intervals cluster around a
characteristic value, giving a high $\rbar$, while successive intervals are
statistically independent. The two cases are distinguished by the \emph{shuffle
test}: randomly permuting the spacings destroys any sequential structure while
preserving the marginal. Genuine repulsion is degraded by shuffling (the rigid
spectrum has obs$<$shuf); a scalar clock is shuffle-invariant (obs$\approx$shuf).
Version~1 never ran this test.

The deeper point is structural. By the Bohigas--Giannoni--Schmit criterion
\citep{bgs1984}, GOE statistics are the fingerprint of a \emph{chaotic,
strongly-mixed} quantum (or wave) spectrum. A classical relaxation oscillator is
integrable, not chaotic; its inter-event times form a renewal process. A
charge-and-release ore reservoir is therefore, on physical grounds, a scalar
renewal clock and \emph{cannot} be a GOE spectrum, whatever the numerical value of
$\rbar$. This erratum makes that argument quantitative on the original data.

\section{Data and methods}

\subsection{Datasets (unchanged from v1)}
We use the same compiled multi-stage absolute ages for four provinces:
(1)~\emph{Andean porphyry copper} (subduction; 11 districts, Re--Os/U--Pb/Ar--Ar)
\citep{maksaev2004,romero2011,deckart2013};
(2)~\emph{South China Nanling W--Sn} (intracontinental granite; 15 deposits)
\citep{mao2017,zhou2016};
(3)~\emph{Tethyan porphyry copper} (collisional; 9 districts across Tibet, Iran,
Pakistan) \citep{hou2009,yang2017};
(4)~\emph{orogenic gold} (metamorphic, non-magmatic; 10 deposits across the
Yilgarn, Abitibi, and West African cratons) \citep{vielreicher2010,pelletier2022,oberthur1998}.
Each district/deposit is treated as one source. These are representative
literature ages, not a systematic census.

\noindent\textbf{What the within-deposit ``intervals'' actually are.} It is
important to be explicit. Within a deposit, the intervals are differences between
ages obtained by \emph{different geochronometers} (U--Pb xenotime, U--Pb
monazite, Ar--Ar muscovite, Re--Os sulfide) on what is often a single
mineralizing episode. For Golden Mile, for example, the four ages
$2675,2655,2645,2630$~Ma give ``intervals'' of $20,10,15$~Myr---comparable to the
individual $\pm4$--$5$~Myr age uncertainties. The quantity version~1 analysed is
therefore dominated by inter-method geochronological scatter, not by resolved
province-scale recurrence.

\subsection{RMT spacing ratio and the diagnostics v1 omitted}
For each source we sort the ages, take inter-event intervals
$s_i=x_{i+1}-x_i$, normalize each source to unit mean, and pool. We compute the
unfolding-free consecutive spacing ratio
\begin{equation}
\rbar=\Big\langle \frac{\min(s_i,s_{i+1})}{\max(s_i,s_{i+1})}\Big\rangle ,
\end{equation}
with reference values Poisson $0.386$, GOE $0.531$, GUE $0.603$
\citep{oganesyan2007,atas2013}. To this we add four diagnostics:
\begin{enumerate}
\item \textbf{Shuffle test.} Randomly permute the pooled spacings (preserving the
marginal, destroying the sequence) and recompute $\rbar$. We report
obs$-$shuf. A scalar clock gives obs$\approx$shuf; genuine GOE gives a small
negative obs$-$shuf (spectral rigidity), resolvable only for $n\gtrsim300$.
\item \textbf{CV and gamma renewal clock.} We compute CV~$=\sigma/\mu$ of the
pooled spacings and compare $\rbar$ to that of a zero-repulsion i.i.d.\ gamma
renewal process of the same CV (shape $k=\mathrm{CV}^{-2}$), sampled at the same
$n$.
\item \textbf{Closed-data control.} We push a zero-repulsion i.i.d.\ gamma null
through the \emph{same} within-deposit mean-normalization ($s_i\to s_i/\bar s$ per
deposit) and pooling, and measure its obs$-$shuf. Mean-normalization forces
$\sum_i s_i=\text{const}$ per deposit, a closed/compositional constraint that
induces negative correlation by construction.
\item \textbf{Age-error Monte Carlo.} We jitter every age within its quoted
$1\sigma$ error (Gaussian), recompute the pooled $\rbar$, and report its
distribution.
\end{enumerate}
All diagnostics are implemented in \texttt{code/audit\_scalar\_clock.py}; the
original pooled $\rbar$ is reproduced by \texttt{code/metallogeny\_rmt.py}.

\section{Results}

\subsection{The spacing ratios are reproduced---and are shuffle-invariant}
We reproduce the version~1 spacing ratios exactly (Table~\ref{tab:audit}). But the
shuffle test, absent from version~1, removes the level-repulsion reading: at every
province obs$\approx$shuf (Figure~\ref{fig:scalar}b). For orogenic gold,
$\rbar_{\rm obs}=0.678$ while $\rbar_{\rm shuf}=0.709$, i.e.\ obs$-$shuf$=-0.031$;
the others are $+0.040$ (Andean), $-0.001$ (Nanling), $+0.074$ (Tethyan). There is
no consistent negative rigidity. The high spacing ratios are a property of the
\emph{marginal} spacing distribution, not of the sequence.

\begin{table}[H]
\centering
\caption{Audit of the four provinces. The version~1 spacing ratios are reproduced
exactly, but are shuffle-invariant (column obs$-$shuf), are reproduced by a
zero-repulsion gamma clock of the same CV, and are not resolvable within age error
(MC range). $n$ is the number of pooled within-source intervals.}
\label{tab:audit}
\small
\begin{tabular}{lcccccc}
\toprule
Province (setting) & $n$ & $\rbar_{\rm obs}$ & obs$-$shuf & CV & gamma clock $\rbar$ & age-error MC \\
 & & (=v1) & (shuffle) & & (95\% band) & $\rbar$ [2.5--97.5\%] \\
\midrule
Andean porphyry Cu (subduction)    & 40 & 0.601 & $+0.040$ & 0.545 & 0.58 [0.49,\,0.67] & [0.37,\,0.56] \\
Nanling W--Sn (intracontinental)   & 36 & 0.574 & $-0.001$ & 0.578 & 0.56 [0.47,\,0.66] & [0.35,\,0.57] \\
Tethyan porphyry Cu (collision)    & 21 & 0.712 & $+0.074$ & 0.509 & 0.60 [0.48,\,0.71] & [0.37,\,0.65] \\
Orogenic gold (\emph{non-magmatic})& 22 & 0.678 & $-0.031$ & 0.310 & 0.73 [0.63,\,0.81] & [0.39,\,0.67] \\
\bottomrule
\end{tabular}
\end{table}

\subsection{A zero-repulsion clock of the same CV reproduces every value}
The spacing ratio is fixed by the CV of the marginal distribution. A zero-repulsion
i.i.d.\ gamma renewal process with the same CV, sampled at the same $n$, brackets
the observed $\rbar$ in all four provinces (Table~\ref{tab:audit};
Figure~\ref{fig:scalar}a): every province sits inside the finite-$n$ 95\% band of
the scalar clock. The orogenic-gold case is the clearest---its low CV ($0.31$)
\emph{is} its high $\rbar$ ($0.68$); no repulsion is needed or implied.

\begin{figure}[H]
\centering
\includegraphics[width=\textwidth]{figs/fig1_scalar_clock.pdf}
\caption{(a)~The spacing ratio is set by the CV alone. The solid curve is $\rbar$
for a zero-repulsion i.i.d.\ gamma renewal clock as a function of CV, with the
finite-$n$ 95\% band shaded; all four provinces sit inside it. (b)~Observed vs
shuffled $\rbar$. obs$\approx$shuf at every province (obs$-$shuf printed above the
bars): the signal is in the marginal distribution, not the sequence.}
\label{fig:scalar}
\end{figure}

\subsection{The apparent ``rigidity'' is a closed-data artifact}
The one province with a notable negative obs$-$shuf is orogenic gold
($-0.031$). This is not spectral rigidity. Pushing a \emph{zero-repulsion} i.i.d.\
gamma null through version~1's own within-deposit mean-normalization and pooling
yields obs$-$shuf$=-0.028\pm0.021$ for the gold configuration---statistically
identical to the observed $-0.031$. Mean-normalizing each deposit's intervals to
unit mean forces their sum to be constant, a closed (compositional) constraint
that induces a spurious negative lag-1 correlation. The negative value is
manufactured by the construction, not by the data.

\subsection{The point estimate is geochronologically unresolvable}
Jittering each age within its quoted $1\sigma$ error and recomputing the pooled
$\rbar$ produces broad distributions (Table~\ref{tab:audit};
Figure~\ref{fig:mc}c). For orogenic gold the central 95\% of the Monte Carlo
spans $\rbar\in[0.39,0.67]$---from the Poisson value to above GUE---with a median
of $0.53$. The version~1 reported value ($0.678$) sits at the very top edge of
that envelope. At $n=21$--$40$, far below the $n\gtrsim300$ needed to separate a
scalar clock from GOE by rigidity, the spacing ratio simply cannot bear the
interpretive weight version~1 placed on it.

There is a deeper, instrument-level reason the sequential reading was never
admissible. Genuine level repulsion is a property of the \emph{ordering} of
closely spaced levels, so detecting it demands an absolute time resolution finer
than the spacings being ordered. The chronometers underlying these compilations---
U--Pb on zircon and xenotime, $^{40}$Ar/$^{39}$Ar on muscovite, Re--Os on
arsenopyrite and molybdenite---carry intrinsic $1\sigma$ uncertainties of order
$\pm 1$--$5$~Ma, which are comparable to or larger than the within-source
intervals themselves. Any micro-scale sequential rigidity that a Wigner--Dyson
spectrum would require is therefore convolved away by the dating errors before it
could be measured: the analytical resolution available in ore geochronology is, on
its face, insufficient to support a level-repulsion statistic. The high $\rbar$
that does survive this smearing is exactly what a narrow marginal (a scalar clock)
produces, because $\rbar$ from a narrow CV is robust to jitter while sequential
rigidity is not.

\subsection{What survives: superposition $\to$ Poisson}
The superposition result is unaffected and we retain it. When intervals from many
independent sources are pooled \emph{without} source isolation, the spacing ratio
collapses toward Poisson (Figure~\ref{fig:mc}d): mixed $\rbar=0.47$ (Andean),
$0.31$ (Nanling), $0.39$ (Tethyan), $0.43$ (gold), against single-source values of
$0.57$--$0.71$. This is the standard superposition theorem---the sum of many
independent renewal processes tends to Poisson \citep{daley2003}---and it is a
genuine, mechanism-independent feature of the data. It says nothing about level
repulsion; it says that mixing erases the marginal clock of each source.

Removing the spurious single-source GOE claim does not weaken this result; it
strengthens it. The superposition collapse never depended on the local statistic
being GOE---only on each source carrying \emph{some} non-Poisson marginal clock,
which the scalar-clock reading supplies just as well. With the pseudo-repulsion
artifact cleared away, the one robust, scale-independent statement the pipeline
supports stands on firmer ground: independent metallogenic sources, superposed,
lose their individual rhythms and the pooled timing relaxes to a memoryless
Poisson process. That is the finding worth carrying forward.

\begin{figure}[H]
\centering
\includegraphics[width=\textwidth]{figs/fig2_age_error_superposition.pdf}
\caption{(c)~Age-error Monte Carlo. Box plots of the pooled $\rbar$ when ages are
jittered within their quoted $1\sigma$ errors; stars mark the version~1 reported
values. The estimates float across the Poisson--GUE range; the v1 gold value sits
at the top edge. (d)~The retained, real result: single (isolated-source) vs mixed
(superposed) $\rbar$. Superposition drives the local statistic toward Poisson
($0.386$).}
\label{fig:mc}
\end{figure}

\section{Discussion}

\subsection{Scalar clock versus level repulsion}
The corrected picture places these systems firmly at the scalar-clock end of the
diagnostic axis. A high $\rbar$ that is shuffle-invariant, CV-determined, and
reproduced by an i.i.d.\ renewal clock is the signature of a narrow marginal
spacing distribution---a system with a characteristic recurrence time and
independent successive intervals. That is precisely what a charge-and-release
reservoir is: a relaxation oscillator. It is the same object as a dripping tap,
a geyser, a glacial j\"okulhlaup cycle, or a nucleosome spacing ruler---all of
which give high $\rbar$ for the same reason and none of which is a GOE spectrum.
Genuine level repulsion, by contrast, requires a chaotic, strongly-mixed spectrum
(open $d/f$ shells, heavy nuclei, hydrogen in a strong magnetic field) per the
Bohigas--Giannoni--Schmit criterion \citep{bgs1984}, and shows resolvable spectral
rigidity (obs$<$shuf) at large $n$. The ore data show none of this.

The distinction is one of dynamical class, not of degree. A charge-and-release
reservoir has a single effective degree of freedom---an accumulating budget that
discharges past a threshold---so its dynamics live on a classical limit cycle and
its inter-event statistics are those of a one-dimensional renewal process.
Wigner--Dyson statistics, on the other hand, are the fingerprint of a
high-dimensional, strongly-coupled, non-integrable system whose levels genuinely
interact. Reading a metallogenic clock as GOE therefore mistakes a classical,
effectively isolated relaxation oscillation for quantum-grade strong coupling---it
assigns a one-degree-of-freedom threshold process to a universality class reserved
for chaotic many-body spectra. The general lesson for applied RMT is that a high
spacing ratio is a necessary but not sufficient condition for level repulsion:
without the shuffle test, a CV check, and a renewal-clock null, any narrow-marginal
recurrence will masquerade as a repelling spectrum.

\subsection{Why the non-magmatic ``control'' does not rescue the claim}
Version~1's strongest argument was that orogenic gold---non-magmatic---shows the
same high $\rbar$, supposedly cutting off an ``inherited from magma'' explanation
and implying intrinsic repulsion. The argument fails because the premise (that the
high $\rbar$ \emph{is} repulsion) is false. Orogenic gold has the same high
$\rbar$ as the magmatic provinces because it is the same kind of object---a scalar
renewal clock (here, a low-CV cluster of multi-method ages on a few mineralizing
episodes), not because it shares an intrinsic repulsion mechanism. The
cross-cratonic consistency of the gold $\rbar$ is consistency of a scalar clock,
not of level repulsion.

\subsection{What can still be said}
Two statements remain defensible. First, individual ore systems behave as scalar
charge-release clocks: their pulse intervals cluster around a characteristic
recurrence time, consistent with reservoir recharge dynamics, but with no evidence
of sequential level repulsion at present resolution. Second, superposition of many
sources randomizes timing toward Poisson, so any clock signature is a single-source
phenomenon. Neither statement invokes random-matrix universality, and the paper no
longer claims it.

\subsection{Limitations}
Beyond the corrected interpretation: sample sizes are small ($n=21$--$40$ pooled
intervals); dating methods are heterogeneous and the within-deposit intervals are
dominated by inter-method scatter; ages are representative literature values, not a
census; and the geochronological errors alone make the spacing ratio unresolvable.
A genuine test of recurrence structure would need many tens of resolved,
single-method pulse ages within one system---currently unavailable.

\section{Conclusions}
We retract the version~1 claim that single ore-forming systems show GOE/GUE level
repulsion. The high spacing ratios are reproduced exactly but are shuffle-invariant,
CV-determined, reproduced by a zero-repulsion renewal clock, contaminated by a
closed-data normalization artifact in their apparent ``rigidity,'' and unresolvable
within age error. A charge-and-release reservoir is a relaxation oscillator---a
scalar renewal clock---which by the Bohigas--Giannoni--Schmit criterion is
structurally not a GOE spectrum. What survives is the superposition result:
pooling independent sources drives the timing toward Poisson. We keep the data and
the superposition analysis, and reframe single ore systems as scalar
charge-release clocks rather than level-repulsive spectra. The correction follows
the scalar-clock audit, Zenodo \url{https://doi.org/10.5281/zenodo.20980670}.

\section*{Code and Data Availability}
All code (including the new \texttt{code/audit\_scalar\_clock.py}), compiled
datasets, and figures: \url{https://github.com/Ruqing1963/metallogeny-rmt}.
Version~1 is preserved in the git history (tag \texttt{v1.0}), in
\texttt{paper/archive/v1/}, and as the prior Zenodo version. Methodological basis
of the correction: Zenodo \url{https://doi.org/10.5281/zenodo.20980670}.

\begin{thebibliography}{20}
\bibitem[At\'as et al.(2013)]{atas2013} At\'as, Y.Y., Bogomolny, E., Giraud, O., Roux, G. (2013). Distribution of the ratio of consecutive level spacings. \textit{Phys. Rev. Lett.}, 110, 084101.
\bibitem[Bohigas et al.(1984)]{bgs1984} Bohigas, O., Giannoni, M.-J., Schmit, C. (1984). Characterization of chaotic quantum spectra and universality of level fluctuation laws. \textit{Phys. Rev. Lett.}, 52, 1--4.
\bibitem[Chelle-Michou et al.(2017)]{chelle2017} Chelle-Michou, C., Rottier, B., Caricchi, L., Simpson, G. (2017). Tempo of magma degassing and porphyry copper genesis. \textit{Scientific Reports}, 7, 40566.
\bibitem[Chen(2026)]{chen2026audit} Chen, R. (2026). A scalar clock is not level repulsion: a fourth spurious-GOE hazard and a re-audit of metallogenic timing. Zenodo, \url{https://doi.org/10.5281/zenodo.20980670}.
\bibitem[Daley \& Vere-Jones(2003)]{daley2003} Daley, D.J., Vere-Jones, D. (2003). \textit{An Introduction to the Theory of Point Processes}, Vol. I. Springer.
\bibitem[Deckart et al.(2013)]{deckart2013} Deckart, K., Clark, A.H., Cuadra, P., Fanning, M. (2013). Time-space evolution of the Rio Blanco-Los Bronces porphyry Cu-Mo cluster. \textit{Mineralium Deposita}, 48, 57--79.
\bibitem[Hou et al.(2009)]{hou2009} Hou, Z., et al. (2009). The Miocene Gangdese porphyry copper belt, southern Tibet. \textit{Ore Geology Reviews}, 36, 25--51.
\bibitem[Maksaev et al.(2004)]{maksaev2004} Maksaev, V., et al. (2004). New chronology for El Teniente. \textit{SEG Special Publication}, 11, 15--54.
\bibitem[Mao et al.(2017)]{mao2017} Mao, J., et al. (2017). Granite-related W-Sn mineralization in South China. \textit{Ore Geology Reviews}, 81, 1045--1058.
\bibitem[Mehta(2004)]{mehta2004} Mehta, M.L. (2004). \textit{Random Matrices}, 3rd ed. Elsevier.
\bibitem[Oberth\"ur et al.(1998)]{oberthur1998} Oberth\"ur, T., et al. (1998). Ages of gold mineralization, Ashanti belt, Ghana. \textit{Economic Geology}, 93, 289--309.
\bibitem[Oganesyan \& Huse(2007)]{oganesyan2007} Oganesyan, V., Huse, D.A. (2007). Localization of interacting fermions at high temperature. \textit{Phys. Rev. B}, 75, 155111.
\bibitem[Pelletier et al.(2022)]{pelletier2022} Pelletier, M., et al. (2022). U--Pb vein xenotime geochronology of orogenic gold, Malartic-Val-d'Or, Abitibi. \textit{Mineralium Deposita}, 57, 1469--1490.
\bibitem[Romero et al.(2011)]{romero2011} Romero, B., et al. (2011). Re-Os geochronology of the Escondida porphyry deposits. \textit{Resource Geology}, 61, 91--100.
\bibitem[Sillitoe(2010)]{sillitoe2010} Sillitoe, R.H. (2010). Porphyry copper systems. \textit{Economic Geology}, 105, 3--41.
\bibitem[Vielreicher et al.(2010)]{vielreicher2010} Vielreicher, N., Groves, D., McNaughton, N. (2010). Timing of gold mineralization, Yilgarn craton. \textit{Mineralium Deposita}, 50, 391--428.
\bibitem[Wigner(1957)]{wigner1957} Wigner, E.P. (1957). Characteristic vectors of bordered matrices. \textit{Ann. Math.}, 62, 548--564.
\bibitem[Yang et al.(2017)]{yang2017} Yang, Z., et al. (2017). Cyclic magmatic-hydrothermal evolution of the Qulong porphyry Cu-Mo deposit, Tibet. \textit{Economic Geology}, 112, 1419--1440.
\bibitem[Zhou et al.(2016)]{zhou2016} Zhou, Y., et al. (2016). Dating of the Xitian W-Sn deposit, South China. \textit{Ore Geology Reviews}, 78, 1--13.
\end{thebibliography}

\end{document}
